\pdfmapfile{+symbols.map}
\pdfmapfile{+euler.map}
\pdfmapfile{+cm.map}
\pdfmapfile{+cmextra.map}
\pdfmapfile{+lm.map}
\documentclass[11pt]{article}
\usepackage[a4paper,margin=24mm]{geometry}
\usepackage[T1]{fontenc}
\usepackage{lmodern,microtype}
\setlength{\emergencystretch}{2em}
\usepackage{amsmath,amssymb,amsthm}
\usepackage[colorlinks=true,urlcolor=blue,linkcolor=blue]{hyperref}
\usepackage{enumitem}
\usepackage{fancyhdr}
\setlength{\parindent}{0pt}
\setlength{\parskip}{6pt}
\setlist{nosep,leftmargin=*}
\pagestyle{fancy}
\fancyhf{}
\fancyhead[L]{\small TLMC 00000000506}
\fancyhead[R]{\small A singleton first Betti support}
\fancyfoot[C]{\thepage}
\setlength{\headheight}{14pt}
\newtheorem{proposition}{Proposition}
\newtheorem{lemma}{Lemma}
\newtheorem{theorem}{Theorem}
\newcommand{\F}{\mathbb F_2}
\newcommand{\m}{\mathfrak m}
\newcommand{\supp}{\operatorname{supp}}
\newcommand{\id}{00000000506}
\begin{document}
\begin{center}
{\LARGE A singleton first Betti support}\par\vspace{5pt}
{\large A disproof of TLMC conjecture \id}\par\vspace{5pt}
{\small Complete self-contained Lean 4 formalization and mathematical report}\par
{\small 21 October 2026; prepared for review, no organizer acceptance claimed}
\end{center}

\section*{1. Claim, convention, and counterexample}
The conjecture asserts, in both its English and Chinese versions, that the
support of the first Betti row of a monomial ideal is connected and has no
isolated points. It adds a proposed explanation involving the Taylor
resolution. We refute the explicit no-isolated-points clause. No interpretation
or verification of the additional Taylor-resolution clause is needed to refute
the conjunction.

We use the standard total grading and the Betti numbers of the ideal as a
module, not the Hilbert function of its entire syzygy module. For a minimal
graded presentation $F_0\twoheadrightarrow I$ with kernel $K$, the first Betti
support is
\[
 \{j\in\mathbb Z:(K/\m K)_j\ne0\}.
\]
This is the minimal-generator definition of the graded Betti row; it agrees
with the shifts in a minimal graded free resolution and with graded Tor
dimensions. See the definitions in \cite{madden,eisenbud}. The formal proof
computes this quotient explicitly rather than assigning a numerical table.

\begin{theorem}
Let $R=\F[X,Y]$, with $\deg X=\deg Y=1$, and let $I=(X,Y)=\m$.
Then $I$ has minimal graded free resolution
\[
 0\longrightarrow R(-2)\xrightarrow{\ d_1\ }
 R(-1)^2\xrightarrow{\ d_0\ }I\longrightarrow0,
 \qquad d_1(h)=(Yh,Xh),\quad d_0(a,b)=Xa+Yb.
\]
Consequently $\beta_{1,2}(I)=1$ and $\beta_{1,j}(I)=0$ for $j\ne2$.
Its first Betti support is the singleton $\{2\}$, which has an isolated point.
\end{theorem}

The field $\F$ is allowed: the source imposes no characteristic restriction.
In characteristic two the usual Koszul relation $(-Y,X)$ is $(Y,X)$.
The ideal is genuinely monomial, generated by the monomials $X$ and $Y$.

\section*{2. Exactness, with every syzygy accounted for}
Represent a polynomial by its finitely supported coefficients
$p_{ij}\in\F$ for $i,j\ge0$. Multiplication by $X$ or $Y$ shifts the respective
index by one, with zero on the new boundary.

Surjectivity of $d_0$ follows from the definition of $(X,Y)$. Also
$d_0d_1(h)=XYh+YXh=0$. To compute the entire kernel, suppose $Xa+Yb=0$.
Comparing coefficients gives
\[
 a_{i,0}=0,\qquad b_{0,j}=0,\qquad
 a_{i,j+1}=b_{i+1,j}\quad(i,j\ge0).
\]
Define $h_{ij}=a_{i,j+1}$. This still has finite support. The displayed
identities yield $a=Yh$ and $b=Xh$. Conversely every such pair is a syzygy.
The coefficient identity $(Yh)_{i,j+1}=h_{ij}$ proves uniqueness of $h$ and
injectivity of $d_1$. Thus the sequence is exact, including at its left end.

\newpage
\section*{3. Minimality and the actual graded quotient}
\subsection*{3.1 The original presentation is minimal}
A polynomial lies in $I$ if and only if its constant coefficient is zero.
Indeed, split the terms with positive $X$ exponent from those with $X$ exponent
zero and positive $Y$ exponent. Thus $p\in I$ precisely when $p=Xa+Yb$.

Similarly,
\[
 p\in\m I=X I+Y I
 \quad\Longleftrightarrow\quad p_{00}=p_{10}=p_{01}=0.
\]
For the reverse implication the same term split gives coefficients $a,b$
whose constant coefficients are $p_{10},p_{01}$, hence $a,b\in I$.
It follows that
\[
 I/\m I\ \xrightarrow{\ \sim\ }\ \F^2,
 \qquad [p]\longmapsto(p_{10},p_{01}).
\]
The classes of $X,Y$ are therefore a basis, both in degree one. This verifies
minimality of the generators used to define $K$, a necessary point when
computing first Betti numbers from a presentation.

\subsection*{3.2 Minimality of the differential}
The map $d_1$ is $R$-linear and additive; $d_0$ is also $R$-linear and
additive. The entries $Y,X$ of $d_1$ lie in $\m$. Thus the image of the free differential
lies in $\m R(-1)^2$, the standard minimal-complex condition. There is no
further differential because $d_1$ is injective.

The grading is not implicit: if $h$ is homogeneous of ordinary degree $r$,
then $(Yh,Xh)$ has degree $r+2$ in $R(-1)^2$. This is exactly the shift
$R(-2)$. The augmentation preserves degree. The ring multiplication is also
proved to respect the standard total grading.

\subsection*{3.3 The quotient $K/\m K$}
The established $R$-linear bijection
$d_1:R(-2)\longrightarrow K$ takes $\m$ onto $\m K$.
In detail, finite sums of arbitrary products of an element of $\m$ with an
element of $K$ equal $XK+YK$, since $\m=(X,Y)$. Under $d_1$, membership in this
submodule is equivalent to $h\in\m$, that is, $h_{00}=0$.

Therefore the map
\[
 K/\m K\longrightarrow\F,\qquad [d_1(h)]\longmapsto h_{00}
\]
is a well-defined additive, $\F$-linear bijection. Its inverse sends
$c\in\F$ to $[d_1(c)]$. In particular, the nonzero class of $d_1(1)$ is a
one-element basis. This is a quotient of the actual polynomial syzygy module,
not a separately invented two-element model.

Now a homogeneous syzygy of degree $j$ can have nonzero quotient class only
when $h_{00}=1$. That constant coefficient contributes degree two, forcing
$j=2$. Conversely $d_1(1)$ is homogeneous of degree two and has nonzero class.
This proves, for every integer $j$, including negative degrees,
\[
 (K/\m K)_j\ne0\quad\Longleftrightarrow\quad j=2.
\]
The one-element basis also gives the stated numerical value
$\beta_{1,2}(I)=1$.

\newpage
\section*{4. Why the singleton disproves the source}
A singleton is connected. The argument does \emph{not} claim otherwise.
It contradicts the separate requirement that the support have no isolated
points.

\begin{proposition}
In every ambient topology, the sole point of a singleton subset is isolated
in that subset. With any irreflexive adjacency relation, a one-vertex support
also has no neighbor.
\end{proposition}
\begin{proof}
The whole ambient space is open, and its intersection with the singleton is
the singleton itself. Hence it is an isolating neighborhood. For adjacency,
any putative neighbor in the support must be the same point, contradicting
irreflexivity.
\end{proof}

Both statements are formalized. The topological theorem quantifies over an
arbitrary open-set predicate, requiring only that the whole space is open;
therefore it applies to any genuine topology. The graph theorem applies in
particular to adjacency of consecutive integer degrees. This avoids relying
on an unstated topological-versus-combinatorial convention in the source.

\subsection*{Alternative indexing: the quotient ring}
Some authors tabulate the Betti numbers of $R/I$ instead of $I$. In that
convention the first relation module is $I$, and the first minimal-generator
quotient is $I/\m I$. Section 3.1 shows that it has two independent generators,
both in degree one. Thus the alternative first support is $\{1\}$, again a
singleton with an isolated point. The Lean source independently proves that
the homogeneous nonzero-residue support of $I/\m I$ is exactly $\{1\}$.

\subsection*{What is not being conflated}
The Hilbert function of the \emph{whole} syzygy module $K\cong R(-2)$ is
different from its Betti row. It is nonzero in every degree $j\ge2$.
The conjecture explicitly names a Betti table, so the standard graded
minimal-generator convention is used. The proof does not replace Betti
support by the Hilbert support of $K$.

We do not need a general construction of Tor, the Taylor resolution, or a
proof of uniqueness of every minimal free resolution. Instead, the exact
minimal graded resolution and the defining minimal-generator quotient are
verified directly for this particular monomial ideal. This is sufficient to
refute the universal claim.

\section*{5. Formal model and scope}
The project uses only Lean 4 core and Std, with no Mathlib dependency.
\begin{itemize}
\item A polynomial is a coefficient function $\mathbb N^2\to\{0,1\}$ with an
existential finite-support bound. Equality ignores the chosen bound.
\item Addition is coefficientwise exclusive-or. Multiplication is obtained
from finite lists of monomials by adding exponents and collecting parity.
Every finitely supported coefficient family has a list representation;
representation independence and all commutative-ring laws are proved.
\item The coefficient operations are proved to satisfy the two-element field
laws. The free rank-two module has the actual polynomial scalar action.
\item The ideal, kernels, finite-sum ideal products, grading, and quotient
equivalence relation are explicit. The final support is not postulated.
\end{itemize}

\newpage
\section*{6. Proof map and reproducibility}
The final theorem is
\texttt{TLMC506.conjecture506\_\allowbreak{}counterexample} in \texttt{Resolution.lean}.
It bundles the monomial-ideal certificate, the verified minimal resolution,
the integer-indexed support identity, and failure of the no-isolated clause.

Key supporting results are:
\begin{itemize}
\item \texttt{polynomial\_\allowbreak{}ring\_\allowbreak{}laws}, \texttt{free\_\allowbreak{}module\_\allowbreak{}laws}, and
\texttt{homogeneous\_\allowbreak{}mul}: actual ring, module, and grading structures.
\item \texttt{kernel\_\allowbreak{}free\_\allowbreak{}rank\_\allowbreak{}one}: unique polynomial parametrization of
every syzygy; \texttt{presentation\_\allowbreak{}minimal\_\allowbreak{}basis}: the basis of $I/\m I$.
\item \texttt{maximal\_\allowbreak{}times\_\allowbreak{}syzygy\_\allowbreak{}is\_\allowbreak{}ideal\_\allowbreak{}product}: identification
with the genuine finite-sum ideal-module product $\m K$.
\item \texttt{quotient\_\allowbreak{}basis}: unique coefficient for every class in $K/\m K$;
\texttt{first\_\allowbreak{}betti\_\allowbreak{}supportZ\_\allowbreak{}iff}: exact support over all integer degrees.
\item \texttt{first\_\allowbreak{}betti\_\allowbreak{}supportZ\_\allowbreak{}graph\_\allowbreak{}isolated}: the graph reading;
\texttt{quotient\_\allowbreak{}first\_\allowbreak{}betti\_\allowbreak{}supportZ\_\allowbreak{}iff}: the alternative indexing.
\end{itemize}

\subsection*{Verification procedure}
The toolchain is pinned to \texttt{leanprover/lean4:v4.31.0}. Run
\texttt{./VERIFY.sh} with that toolchain available on \texttt{PATH}, or set
\texttt{LEAN\_\allowbreak{}BIN\_\allowbreak{}DIR} to its \texttt{bin} directory. The script copies
\emph{only authored sources and project configuration} into a fresh temporary
directory, then performs:
\begin{enumerate}
\item a full Lake build and strict, warning-as-error source checks;
\item an environment-wide transitive axiom audit of every authored theorem,
including generated supporting declarations;
\item official \texttt{leanchecker Main} kernel replay of the fresh build;
\item independent finite polynomial tests in \texttt{reproduce.py}.
\end{enumerate}
The permitted axiom set is exactly Lean's standard
\texttt{propext}, \texttt{Classical.choice}, and \texttt{Quot.sound}.
There are no added axioms, unfinished proofs, native evaluation proof escapes,
or kernel bypasses. The Python checks are supplementary; no finite experiment
is used to establish the universal theorem. Verification logs and a SHA-256
manifest are included in the package.

\newpage
\subsection*{Repository eligibility and provenance}
The dated read-only eligibility check used upstream commit
\texttt{b355e42906e6eaf52bc7e\allowbreak{}9d005e0ef02d8aa7400}.
The exact metadata row for \id{} has both solved flags set to
\texttt{false}. The source blob is
\texttt{6a9833bfdd4546f0e506\allowbreak{}5e9a5bdaad0680da0c14}.
The solution directory was absent, and padded, unpadded, semantic, and complete
saved PR title/body checks found no relevant prior submission for this ID. A
dated public eligibility record is included. This is a dated repository-status check, not a claim of worldwide
mathematical novelty or a guarantee about later submissions. The elementary
Koszul counterexample uses standard algebra.

\begin{thebibliography}{9}
\bibitem{tlmc} The Last Math Competition,
\href{https://github.com/The-Last-Math-Competition/The-Last-Math-Competition/blob/b355e42906e6eaf52bc7e9d005e0ef02d8aa7400/conjectures/00000000506.md}{Conjecture \id},
English and Chinese source, accessed 21 October 2026.
\bibitem{madden} J. Madden, \emph{Lecture 7: Betti Numbers}, 2006, pp. 1--2.
\href{https://www.math.lsu.edu/~madden/M7290Fall2006/lecture7.pdf}{Primary lecture notes}.
\bibitem{eisenbud} D. Eisenbud, \emph{The Geometry of Syzygies}, Springer, 2005,
Section 1B, especially minimal resolutions and Corollary 1.5.
\href{https://www-users.cse.umn.edu/~reiner/REU/REU2019notes/2005_Book_TheGeometryOfSyzygies.pdf}{Author's book, hosted course copy}.
\end{thebibliography}
\end{document}
